\documentclass[a4paper,12pt]{article}
\usepackage{docsTemplate}
\usepackage{usecase}
\usepackage{float}

\setTitle{Analisi dei Requisiti}
\setAuthors{Filippo Compagno \\ & Sara Gioia Fichera}
\setVerificators{Andrea Masiero \\ & Francesco Balestro}
\setApprovation{}
\setVersion{0.3}
\setType{Esterno}
\setDestination{Tullio Vardanega \\ & Riccardo Cardin \\ & Sanmarco Informatica \\[-0.1em] & Team Three Technologies}

\begin{document}
    \include{docTitlepage}

    \section*{Registro delle modifiche} {
        \begin{table}[h!]
            \rowcolors{2}{lightgray}{white}
                \begin{tabularx}{\textwidth}{|l|l|l|l|L|L|}
                \hline
                \textbf{Vers.} & \textbf{Data} & \textbf{Autore} & \textbf{Ruolo} & \textbf{Verifica} & \textbf{Oggetto} \\
                \hline
                \textbf{0.3} & 24/11/25 & Sara Gioia Fichera & Analista & Francesco Balestro &  Definizione Attori \\
                \hline
                \textbf{0.2} & 20/11/25 & Filippo Compagno & Analista & Francesco Balestro & Caratteristiche degli utenti \\
                \hline
                \textbf{0.1} & 07/11/25 & Sara Gioia Fichera & Analista & Andrea Masiero & Struttura iniziale del documento \\
                \hline
            \end{tabularx}
        \end{table}
    }
    \newpage
    
    \tableofcontents
    \newpage

    \listoffigures
    \newpage
    
    \listoftables
    \newpage

    \section{Introduzione} {
        \subsection{Scopo del documento} {        
            Questo documento ha lo scopo di descrivere in maniera esaustivo e precisa i requisiti funzionali e non funzionali del sistema software da sviluppare per il progetto DIPReader. Per questo motivo integra un'attenta descizione dei casi d'uso, che sono fonte di requisiti.
            Tale analisi è il risultato di un attento studio del capitolato, integrato con quanto concordato durante gli incontri svolti con Sanmarco Informatica, la proponente.
        }
        \subsection{Glossario} {
            Si rimanda al documento Glossario per chiarire i termini nel dominio del progetto che possono risultare ambigui o di difficile comprensione. I vocaboli presenti nel glossario sono stati segnalati seguendo la notazione:
            \begin{center}
                parola\textsubscript{G}
            \end{center}
        }
        \subsection{Riferimenti} {
            \subsubsection{Riferimenti normativi} {
                \begin{itemize}
                    \item Norme di Progetto
                    \item \hypertarget{capitolato}{Capitolato}
                    \begin{itemize}
                        \item {\url{https://www.math.unipd.it/~tullio/IS-1/2025/Progetto/C3.pdf}}
                        \item[] (Ultimo accesso: 2025/11/07)
                    \end{itemize}
                    \item Regolamento del Progetto Didattico
                    \begin{itemize}
                        \item \url{https://www.math.unipd.it/~tullio/IS-1/2025/Dispense/PD1.pdf}
                        \item[] (Ultimo accesso: 2025/11/07)
                    \end{itemize}
                    \item 830-1998 - IEEE Recommended Practice for Software Requirements Specifications
                    \begin{itemize}
                        \item \url{https://ieeexplore.ieee.org/document/720574/references#references}
                        \item[] (Ultimo accesso: )
                    \end{itemize}
                \end{itemize}
            }
            \subsubsection{Riferimenti informativi} {
                \begin{itemize}
                    %\item Slide del corso: T4 - Gestione di progetto
                    %\begin{itemize}
                    %    \item \url{https://www.math.unipd.it/~tullio/IS-1/2025/Dispense/T04.pdf}
                    %    \item[] (Ultimo accesso: 2025/11/07)
                    %\end{itemize}
                    \item Slide del corso: T5 - Analisi dei requisiti
                    \begin{itemize}
                        \item \url{https://www.math.unipd.it/~tullio/IS-1/2025/Dispense/T05.pdf}
                        \item[] (Ultimo accesso: 2025/11/07)
                    \end{itemize}
                    \item Slide del corso: P2 - Diagrammi dei Casi d'Uso
                    \begin{itemize}
                        \item \url{https://www.math.unipd.it/~rcardin/swea/2022/Diagrammi\%20Use\%20Case.pdf}
                        \item[] (Ultimo accesso: 2025/11/07)
                    \end{itemize}
                    \item Repository Github del Prof. Cardin
                    \begin{itemize}
                        \item \url{https://github.com/rcardin/swe-imdb}
                        \item[] (Ultimo accesso: )
                    \end{itemize}
                    \item Glossario
                \end{itemize}
            }
        }
    }

    \newpage
    \section{Descrizione del prodotto} {
        \subsection{Presentazione del prodotto} {
            DIPReader ha lo scopo di supportare la gestione e la conservazione digitale dei documenti, consentendo una transizione strutturata da archivi cartacei a soluzioni cloud avanzate. Il sistema centralizza i documenti e ne assicura autenticità, integrità, leggibilità, reperibilità e sicurezza, offrendo funzionalità di acquisizione, estrazione dei dati e consultazione. Inoltre, permette di richiedere la disponibilità offline dei documenti, garantendo sempre la conformità ai requisiti normativi. 
        }
        \subsection{Funzioni del prodotto} {
            Le funzioni del prodotto sono: 
            \begin{itemize}
                \item Rappresentazione semplificata del contenuto 
                \item Ricerca dei documenti tramite metadati
                \item Visualizzazione in anteprima di alcuni formati
                \item Salvare documenti nel computer dell'Utente
                \item Funzionamento offline
            \end{itemize}
        }
        \subsection{Caratteristiche dell'utente} {
            L’applicazione è destinata a un pubblico ampio e variegato, comprendente ricercatori, professionisti, cittadini interessati o chiunque abbia bisogno di consultare documenti archiviati. Gli utenti possono avere competenze tecniche e digitali molto diverse: dai non esperti, che desiderano trovare rapidamente informazioni, ai professionisti abituati a interpretare metadati complessi. \\
            \\
            Tutti gli utenti condividono l’esigenza di effettuare ricerche efficaci e mirate, utilizzando filtri sui metadati, consultare il contenuto dei documenti e scaricare informazioni in modo affidabile. L’interfaccia dovrà quindi garantire usabilità, chiarezza dei metadati principali e supportare la consultazione e l’estrazione dei documenti in maniera intuitiva.
        }

    }
    \newpage
    \section{Casi d'uso} {
        \subsection{Scopo} {
            Lo scopo di questa sezione è spiegare singolarmente ciascun Caso d'uso individuato.  
        }
        \subsection{Attori} {
            \begin{itemize}
                \item Utente: in questo contesto è l'Attore principale, cioè colui che intergisce attivamente con il sistema e svolge le azioni indicate dal Caso d'Uso.
                \item Motore semantico: si occupa di elaborare una ricerca semantica fatta dall'utente per produrre un risultato. In questo contesto è un Attore secondario, in quanto interagisce passivamente con il sistema.
                \item Gestore di stampa SO: si occupa di gestire la richiesta di stampa di un documento da parte dell'utente. In questo contesto è un Attore secondario, in quanto interagisce passivamente con il sistema.
            \end{itemize}
        }
        \subsection{Casi d'uso} {

        }
    }
    \newpage
    \section{Requisiti} {
       
        \subsection{Requisiti Funzionali} {
            
        }
        \subsection{Requisiti di Qualità} {
           
        }
        \subsection{Requisiti di Vincolo} {
           
        }
    }

\end{document}